%%=============================================================================
%% Voorwoord
%%=============================================================================


%% TODO:
%% Het voorwoord is het enige deel van de bachelorproef waar je vanuit je
%% eigen standpunt (``ik-vorm'') mag schrijven. Je kan hier bv. motiveren
%% waarom jij het onderwerp wil bespreken.
%% Vergeet ook niet te bedanken wie je geholpen/gesteund/... heeft

\chapter*{\IfLanguageName{dutch}{Woord vooraf}{Preface}}%
\label{ch:voorwoord}

Tijdens het uitwerken van deze bachelorproef heb ik ondervonden dat het ontwikkelen van een machine learning model slechts één onderdeel is van een veel breder proces. Waar ik initieel vooral gefocust was op het bouwen en optimaliseren van modellen, werd al snel duidelijk dat het correct formuleren van het probleem en het interpreteren van resultaten minstens even belangrijk zijn.
\newline
\newline
Een van de grootste uitdagingen was het afbakenen van de scope. In het begin wilde ik verschillende pistes combineren, waaronder meerdere datasets, chaining-technieken en transfermogelijkheden. Doorheen het proces heb ik moeten leren om keuzes te maken en te focussen op één duidelijke rode draad. Dit heeft niet alleen de kwaliteit van het werk verbeterd, maar ook mijn inzicht in het onderzoeksproces versterkt.
\newline
\newline
Daarnaast bleek de realiteit van cybersecuritydata complexer dan verwacht. Vooral bij datasets met zwakke labels en sterke class imbalance werd duidelijk dat klassieke evaluatiemetrics niet altijd een correct beeld geven van de performantie van een model. Dit dwong mij om verder te kijken dan enkel accuracy en F1-scores, en om meer aandacht te besteden aan de praktische waarde van resultaten, zoals het correct prioritiseren van events en het reconstrueren van verbanden.
\newline
\newline
Tijdens het project heb ik ook verschillende technische obstakels moeten overwinnen. Problemen zoals foutieve datarepresentaties, onverwachte modeloutput of configuratiefouten zorgden regelmatig voor vertraging. In plaats van deze enkel op te lossen door parameters aan te passen, heb ik geleerd om telkens een stap terug te nemen en mijn aannames kritisch te evalueren. Dit proces van iteratie en debugging heeft een belangrijke rol gespeeld in mijn leertraject.
\newline
\newline
Een belangrijk inzicht dat ik heb meegenomen, is dat goede machine learning niet enkel draait om het behalen van sterke resultaten, maar vooral om het correct interpreteren ervan. Het streven naar hoge metrics maakte geleidelijk plaats voor een meer kritische houding, waarbij de vraag centraal stond wat deze cijfers werkelijk betekenen binnen een realistische context.
\newline
\newline
Tot slot wil ik mijn promotor en begeleiders bedanken voor hun ondersteuning en feedback tijdens dit traject. Hun input heeft mij geholpen om mijn werk kritisch te bekijken en verder te verfijnen. Daarnaast wil ik ook mijn omgeving bedanken voor de steun en motivatie gedurende deze periode.
\newline
\newline
Deze bachelorproef heeft mij niet alleen technische kennis bijgebracht, maar ook geleerd om om te gaan met onzekerheid, iteratief te werken en onderbouwde conclusies te formuleren. Deze ervaring vormt een waardevolle basis voor mijn verdere ontwikkeling binnen het vakgebied.